\section{Random Variables}

\begin{definition}
    A \textbf{random variable } is a map $X: \Omega \to \mathbb{R}$ where $\Omega$ is the set of all events of
    a probability space.
\end{definition}

\begin{example}
    The event that the random variable X takes a value less than or equal to 5 is denoted as:
    $$
        X \leq 5 \coloneqq \{\omega \in \Omega \mid X(\omega) \leq 5\}
    $$
\end{example}

\subsection{Expectation}

\begin{definition}
    Let $X$ be a random variable. The expectation of $X$ is defined as:
    $$
    \mathbb{E}[X] \coloneqq \sum_{\alpha \in W_x} \alpha P(X=\alpha)
    $$
    aslong as the sum is absolutely convergent, else we say the expectation is undefined.
\end{definition}

\begin{theorem}
    Let $X$ be a random variable with values in $\mathbb{N}_0$. Then:
    $$
        \mathbb{E}[X] = \sum_{k=1}^\infty P(X \geq k)
    $$
\end{theorem}

\begin{proof}
    $$
        \mathbb{E}[X] = \sum_{k=1}^\infty k P(X=k) = \sum_{k=1}^\infty \sum_{j=1}^k P(X=k) = \sum_{j=1}^\infty \sum_{k=j}^\infty P(X=k) = \sum_{j=1}^\infty P(X \geq j)
    $$
\end{proof}

\begin{example}
    We throw a coin and we know $P[\text{head}] = p$. We want to know $\mathbb{E}[X]$ with $X \coloneqq \text{number of throws until we get head}$. Using the above theorem we get:
    $$
        \mathbb{E}[X] = \sum_{k=1}^\infty P(X \geq k) = \sum_{k=1}^\infty (1-p)^{k-1} = \frac{1}{1-(1-p)} = \frac{1}{p}
    $$
\end{example}

\begin{definition}
    For an event $E \subseteq \Omega$ the corresponding indicator-variable $X_E$ is defined as:
    $$
        X_E(\omega) = \begin{cases} 1 & \text{if } \omega \in E \\ 0 & \text{if } \omega \notin E \end{cases}
    $$
\end{definition}

\begin{remark}
    For an indicator variable $X_E$ we have:
    $$
        \mathbb{E}[X_E] = P(E)
    $$
\end{remark}

\begin{theorem}
    For random variables $X_1, \dots, X_n$ we have:
    $$
        X \coloneqq \sum_{i=1}^n a_i X_i \mid a_1, \dots, a_n \in \mathbb{R}
    $$
    $$
        \mathbb{E}[X] = \sum_{i=1}^n a_i \mathbb{E}[X_i]
    $$
    This is called the linearity of expectation.
\end{theorem}

\begin{proof}
    $$
        \mathbb{E}[X] = \sum_{\omega \in \Omega} X(\omega) P(\omega) = \sum_{\omega \in \Omega} (a_1 X_1(\omega) + \dots + a_n X_n(\omega)) P(\omega)
    $$
    $$
        = a_1 \sum_{\omega \in \Omega} X_1(\omega) P(\omega) + \dots + a_n \sum_{\omega \in \Omega} X_n(\omega) P(\omega) = a_1 \mathbb{E}[X_1] + \dots + a_n \mathbb{E}[X_n]
    $$
\end{proof}

\begin{example}
    We throw a coin $m$ times. Let $X \coloneqq$ number of partial sequences with threetimes heads. What is the expectation of $X$?
    $$
        X = X_1 + X_2 + \dots + X_{m-2}
    $$
    where $X_i$ is the indicator variable for the event that the $i$-th partial sequence is three heads.
    $$
        \mathbb{E}[X_i] = \frac{1}{8}
    $$
    By the linearity of the expectation we get:
    $$
        \mathbb{E}[X] = \sum_{i=1}^{m-2} \mathbb{E}[X_i] = \frac{m-2}{8}
    $$
\end{example}

\begin{algorithm}
#Recall Stable sets in a graph is a set of vertices that are not adjacent to each other. We will now use a randomized algorithm to determin a maximal stable set.

S[v] = 1/0 with probability p
for all u,v in E:
    if S[u] = 1 and S[v] = 1:
        S[v] = 0
return S
\end{algorithm}

this results in a stable set but what is the expected size of S?

\begin{theorem}
    The algorithm produces a stable set of expected size at least $n*P-m*P^2$
\end{theorem}

\begin{proof}
    Let $X$ be the random variable for the size of the stable set $S$ after the first step. Further let $Y$ be the random variable for the amount of edges we looked at in step 2. Observe
    $$
\mid S \mid  \geq X - Y
    $$
    By the linearity of expectation we get:
    $$
        \mathbb{E}[\mid S \mid] \geq \mathbb{E}[X] - \mathbb{E}[Y]
    $$
    We know that $\mathbb{E}[X] = nP$. Further we know that $\mathbb{E}[Y] = mP^2$. Thus we get:
    $$
        \mathbb{E}[\mid S \mid] \geq nP - mP^2
    $$
\end{proof}

\subsection{Variance}

\begin{definition}
    For a random variable $X$ with $\mathbb{E}[X] = \mu$ the variance is defined as:
    $$
        Var[X] \coloneqq \mathbb{E}[(X-\mu)^2]
    $$
    The standard deviation is defined as $\sigma \coloneqq \sqrt{Var[X]}$.
\end{definition}

\begin{theorem}
    For a random variable $X$ we have $Var[X] = \mathbb{E}[X^2] - \mathbb{E}[X]^2$
\end{theorem}

\begin{lemma}
    For a random variable $X$ and $a, b \in \mathbb{R}$ we have $Var[aX+b] = a^2 Var[X]$
\end{lemma}

\subsection{Multiple Random Variables}

\begin{example}
    Throw a coin and a dice. Let $X$ be the outcome of the coin and $Y$ be the outcome of the dice. What is the expectation of $X+Y$?
\end{example}

\begin{definition}
    Random variables $X_1, X_2, \dots, X_n$ are called independent iff for all $a_1, \dots, a_n \in W_{X_1} \times \dots \times W_{X_n}$ we have:
    $$P[X_1=a_1, \dots, X_n=a_n] = P[X_1=a_1] \dots P[X_n=a_n]$$
\end{definition}

Observe that events $A_1, \dots, A_n$ are independent iff their indicator variables $X_{A_1}, \dots, X_{A_n}$ are independent.

\begin{lemma}
    Let $X_1, \dots, X_n$ be independent random variables and $S_1, \dots, S_n \subseteq \mathbb{R}$ be arbitrary sets then:
    $$P[X_1 \in S_1, \dots, X_n \in S_n] = P[X_1 \in S_1] \dots P[X_n \in S_n]$$
    Hence for $X_1, \dots, X_n$ independent random variables and $I \subseteq \{1, \dots, n\}$ we have $(X_i)_{i \in I}$ are independent.
\end{lemma}

\begin{proof}
    
\end{proof}

\begin{theorem}
    Let $f_1, \dots, f_n: \mathbb{R} \to \mathbb{R}$ be functions in $\mathbb{R}$ and $X_1, \dots, X_n$ be independent random variables. Then $f_1(X_1), \dots, f_n(X_n)$ are independent too.
\end{theorem}

\begin{proof}
\end{proof}
    
\subsection{Combined Random Variables}

\begin{theorem}
    For two independent random variable $X, Y$ and $Z = X+Y$ we have:
    $$f_Z(\alpha) = \sum_{\beta \in W_X} f_X(\beta) f_Y(\alpha-\beta)$$
\end{theorem}

    This can intuitivaly be seen by looking at every possibility for $X = \beta$: for $X+Y = \alpha$ to hold we need $Y = \alpha - \beta$. Since $X$ and $Y$ are independent we can multiply their probabilities and sum that over all $\beta$
    
\begin{theorem}
    For random variables $X_1, \dots , X_n$ we have:
    $$\mathbb{E}[X_1 + \dots + X_n] = \mathbb{E}[X_1] + \dots + \mathbb{E}[X_n]$$
\end{theorem}

\begin{theorem}
    For independent random variablex $X_1, \dots , X_n$ we have:
    $$\mathbb{E}[X_1 \dots X_n] = \mathbb{E}[X_1] \dots \mathbb{E}[X_n]$$
    $$Var[X_1 + \dots + X_n] = Var[X_1] + \dots + Var[X_n]$$
\end{theorem}

\begin{important}
    Note that $Var[X \cdot Y] = Var[X] \cdot Var[Y]$ does not hold in general.
\end{important}